\begin{exercise}{122}
Let $p\colon\mathcal E\to\mathcal C$ be a cocartesian fibration and $K$ a marked simplicial set
equipped with a map $f\colon K\to\mathcal C^{\sharp}$. Consider the sub-$\infty$-category of $\mathrm{Fun}^{\mathrm{mcc}}_f(K,\mathcal E)$ on those
$1$-simplices whose corresponding map $K\times\Delta^1\to\mathcal E$ is a map of marked simplicial sets
$K\times(\Delta^1)^{\sharp}\to\mathcal E^{\natural}$, i.e., we consider only those transformations of functors whose components
are pointwise $p$-cocartesian. Show that this sub-$\infty$-category is given by $\mathrm{Fun}^{\mathrm{mcc}}_f(K,\mathcal E)^{\simeq}$.

\begin{proof}
First, for any cocartesian fibration $q\colon T\to B$, an edge over
$s_0(b)$ is $q$-cocartesian if and only if it is an equivalence in the fibre $T_b$.
The forward implication follows by pulling back to $b$ and applying
Exercises~\ref{ex:117} and~\ref{ex:115}.
Conversely, Exercise~\ref{ex:115} makes a fibre equivalence cocartesian over the point,
and Exercise~\ref{ex:118} makes its image locally $q$-cocartesian.
Since $q$ supplies a cocartesian lift with the prescribed source,
Lemma~3.2.6(2) makes that edge $q$-cocartesian.
This argument imposes no $\infty$-category hypothesis on $B$.

Let $L$ be the underlying simplicial set of $K$, and put
\[
 P=p_*\colon\mathcal E^L\to\mathcal C^L,
 \qquad M=P^{-1}(f),\qquad A=\mathrm{Fun}^{\mathrm{mcc}}_f(K,\mathcal E).
\]
By Proposition~3.2.13 and pullback, $M$ is an $\infty$-category.
Definition~3.2.14 makes $A$ a full subcategory of $M$, so it is an
$\infty$-category and reflects equivalences (Corollary~1.2.78).
For an edge $\alpha$ of $A$, the preceding criterion for $P$, followed by
Proposition~3.2.13, says that $\alpha$ is an equivalence of $A$ precisely
when every component $\alpha_k$ is $p$-cocartesian.

The marked-map condition implies this component condition because every
degenerate edge of $K$ is marked.
Conversely, suppose all components are $p$-cocartesian.
For a marked edge $e\colon k\to k'$ of $K$, the endpoint functors send $e$
to $p$-cocartesian edges.
The triangle in $e\times\Delta^1$ with vertices
$(k,0),(k',0),(k',1)$ therefore has both successive edges $p$-cocartesian.
Exercise~\ref{ex:119} makes its diagonal locally $p$-cocartesian,
and Lemma~3.2.6(2) makes it $p$-cocartesian.
Thus $\alpha$ carries every marked edge of $K\times(\Delta^1)^{\sharp}$
to a marked edge of $\mathcal E^{\natural}$.

The selected edges are consequently exactly the equivalences of $A$.
By Lemma~1.2.27, the sub-$\infty$-category they determine is $A^{\simeq}$,
as asserted.
\end{proof}
\end{exercise}
